% Machine learning terms, in alphabetical order.
\domain{Machine Learning}

\begin{entry}{datasets}{Training, Validation, and Test Data}
  \defn{iso-22989}{\emph{Training data} are samples used to train a machine
    learning model. \emph{Validation data} are samples used to tune
    hyperparameters or compare candidate models. \emph{Test data} are
    samples, independent of the training data, used to assess the performance
    of the final model.}
  \defn[§7.2]{hastie-2009}{In a data-rich setting, the data are divided into
    three parts. The training set fits the models, the validation set
    estimates prediction error for model selection, and the test set assesses
    the generalisation error of the chosen model. The test set should be kept
    \enquote{in a vault} until the end.}
  \defn[§5.3]{goodfellow-2016}{The test set estimates generalisation error
    and must not influence model choices. A validation set, held out from the
    training data, is used to set hyperparameters.}
  \note{In some communities, especially medicine, \enquote{validation} means
    \emph{external} evaluation on new data, which is close to what ML calls
    testing. Check which meaning a document uses.}
  \related{generalization,overfitting,hyperparameter}
\end{entry}

\begin{entry}{deep-learning}{Deep Learning}[DL]
  \defn{lecun-2015}{Representation-learning methods with multiple levels of
    representation, obtained by composing simple non-linear modules that
    each transform the representation at one level into a more abstract
    representation at the next.}
  \defn[ch.~1]{goodfellow-2016}{An approach to AI in which computers learn
    from experience and understand the world as a hierarchy of concepts,
    each defined in relation to simpler ones. The depth of that hierarchy
    gives the field its name.}
  \defn{iso-22989}{An approach to creating rich hierarchical representations
    by training neural networks with many hidden layers.}
  \related{neural-network,machine-learning,transformer}
\end{entry}

\begin{entry}{distribution-shift}{Distribution Shift}[dataset shift, drift]
  \defn{quinonero-2009}{The situation in which the joint distribution of
    inputs and outputs differs between the training and test (or deployment)
    stages. Special cases include covariate shift, where input distributions
    change, and prior probability shift.}
  \defn{gama-2014}{\emph{Concept drift}: a change over time in the
    relationship between input data and the target variable, which makes a
    model learned on past data less accurate.}
  \defn{nist-ai-rmf}{A risk specific to AI: data used to build a system may
    not represent the context of use, and the conditions of deployment may
    change, so that performance and trustworthiness degrade after
    deployment.}
  \note{\enquote{Data drift} usually means a change in the input
    distribution. \enquote{Concept drift} means a change in the input–output
    relationship. Both motivate post-deployment monitoring.}
  \related{generalization,robustness,validity-reliability}
\end{entry}

\begin{entry}{embedding}{Embedding}
  \defn{bengio-2003}{A learned, distributed feature vector that represents a
    discrete symbol, such as a word, in a continuous space, so that similar
    symbols have similar vectors.}
  \defn{mikolov-2013}{A continuous vector representation of words, learned
    from large corpora, in which semantic and syntactic regularities appear
    as vector offsets.}
  \defn[ch.~15]{goodfellow-2016}{Any learned representation that maps
    inputs to points in a vector space where distances reflect a notion of
    similarity useful for a task.}
  \related{deep-learning,retrieval-augmented-generation}
\end{entry}

\begin{entry}{fine-tuning}{Fine-Tuning}
  \defn{howard-ruder-2018}{Adapting a model pre-trained on a large general
    corpus to a target task by continuing to train some or all of its
    parameters on task-specific data.}
  \defn{bommasani-2021}{One form of \emph{adaptation}, in which a foundation
    model's parameters are updated on downstream data. It contrasts with
    prompting, which leaves parameters unchanged.}
  \defn{iso-22989}{Additional training of a pre-trained model on a smaller,
    task- or domain-specific dataset.}
  \related{transfer-learning,foundation-model,rlhf}
\end{entry}

\begin{entry}{generalization}{Generalization}
  \defn[§5.2]{goodfellow-2016}{The ability of a model to perform well on
    previously unobserved inputs, not only on the inputs it was trained on.}
  \defn[ch.~7]{hastie-2009}{Measured by the \emph{generalisation} (test)
    error: the expected prediction error of a model on an independent sample
    from the same distribution as the training data.}
  \defn{iso-22989}{The ability of a machine learning model to make
    predictions that are as accurate on new data as on training data.}
  \note{Classical generalisation assumes that test data come from the same
    distribution as training data. Generalisation under distribution shift is
    a harder, separate problem.}
  \related{overfitting,distribution-shift,datasets}
\end{entry}

\begin{entry}{gradient-descent}{Gradient Descent}
  \defn[§4.3]{goodfellow-2016}{An optimisation method that decreases a
    function by repeatedly moving its parameters a small step in the
    direction of the negative gradient.}
  \defn[§5.2.4]{bishop-2006}{An iterative method for training models in which
    each update moves the weights in the direction of steepest decrease of
    the error function. \emph{Stochastic} or \emph{online} gradient descent
    updates on one data point, or a small batch, at a time.}
  \related{loss-function,neural-network}
\end{entry}

\begin{entry}{hyperparameter}{Hyperparameter}
  \defn[§5.3]{goodfellow-2016}{A setting that controls a learning
    algorithm's behaviour and is not adapted by the algorithm itself, such as
    model capacity or learning rate. Hyperparameters are usually chosen with
    a validation set.}
  \defn{iso-22989}{A characteristic of a machine learning algorithm that
    affects its learning process. It is selected before training rather than
    learned from the training data.}
  \defn{bishop-2006}{A parameter of a prior distribution or of the model
    family, rather than a parameter fitted directly to the data. Examples
    are regularisation coefficients and the number of basis functions.}
  \related{datasets,model}
\end{entry}

\begin{entry}{loss-function}{Loss Function}[cost function, objective]
  \defn[§4.3]{goodfellow-2016}{The function to be minimised during training,
    also called the cost, error, or objective function. It measures how far a
    model's outputs are from the targets.}
  \defn[§2.4]{hastie-2009}{A function $L(Y, f(X))$ that penalises errors in
    prediction. Examples are squared-error loss for regression and 0–1 loss
    for classification.}
  \defn{sutton-barto-2018}{In reinforcement learning, the role of the loss is
    played by a \emph{reward signal} to be maximised rather than an error to
    be minimised.}
  \related{gradient-descent,alignment}
\end{entry}

\begin{entry}{machine-learning}{Machine Learning}[ML]
  \defn*{mitchell-1997}{A computer program is said to learn from experience
    $E$ with respect to some class of tasks $T$ and performance measure $P$,
    if its performance at tasks in $T$, as measured by $P$, improves with
    experience $E$.}
  \defn{samuel-1959}{The field of study that gives computers the ability to
    learn without being explicitly programmed. This wording is widely
    attributed to Samuel but does not appear verbatim in the paper.}
  \defn{iso-22989}{The process of optimising model parameters through
    computational techniques so that the model's behaviour reflects the data
    or experience.}
  \defn[ch.~1]{goodfellow-2016}{The ability of AI systems to acquire their
    own knowledge by extracting patterns from raw data.}
  \note{Mitchell's definition is operational and testable. ISO/IEC 22989
    describes the mechanism, parameter optimisation, and is the one most
    often used in governance documents.}
  \related{artificial-intelligence,supervised-learning,unsupervised-learning,reinforcement-learning}
\end{entry}

\begin{entry}{model}{Model}[ML model]
  \defn{iso-22989}{A mathematical construct that generates an inference or
    prediction from input data. A \emph{machine learning model} is a model
    produced by training on data.}
  \defn[Recital~97]{eu-ai-act}{Distinguished from an \emph{AI system}: a
    model is a component that needs further elements, such as a user
    interface, to become a system.}
  \defn[ch.~1]{bishop-2006}{A parametric function, or a probability
    distribution, whose parameters are fitted to data so that it can predict
    target values for new inputs.}
  \related{ai-system,hyperparameter,model-card}
\end{entry}

\begin{entry}{neural-network}{Neural Network}[artificial neural network]
  \defn{iso-22989}{A network of one or more layers of neurons, connected by
    weighted links with adjustable weights, that takes input data and
    produces an output.}
  \defn[ch.~6]{goodfellow-2016}{A \emph{feedforward network} defines a
    mapping $y = f(x; \theta)$ as a composition of many simpler functions, or
    layers, and learns the parameters $\theta$ that give the best function
    approximation.}
  \defn[ch.~5]{bishop-2006}{A flexible class of parametric non-linear
    functions made of successive linear combinations and non-linear
    activation functions, with parameters fitted by optimisation.}
  \related{deep-learning,transformer,gradient-descent}
\end{entry}

\begin{entry}{overfitting}{Overfitting}
  \defn[§5.2]{goodfellow-2016}{Occurs when the gap between the training
    error and the test error is too large, because the model has memorised
    properties of the training set that do not hold for new data.}
  \defn{iso-22989}{Creating a model that fits the training data too
    precisely and therefore fails to generalise to new data.}
  \defn[§1.1]{bishop-2006}{The phenomenon in which an overly flexible model
    fits noise in the training data, giving small training error but poor
    predictions on new inputs.}
  \related{generalization,datasets}
\end{entry}

\begin{entry}{reinforcement-learning}{Reinforcement Learning}[RL]
  \defn*{sutton-barto-2018}{Reinforcement learning is learning what to
    do---how to map situations to actions---so as to maximize a numerical
    reward signal.}
  \defn{iso-22989}{Learning an optimal sequence of actions to maximise a
    cumulative reward, through trial-and-error interaction with an
    environment.}
  \defn{russell-norvig-2021}{Learning in which an agent learns from a
    series of reinforcements, rewards and punishments, how to act so as to
    maximise expected utility.}
  \related{rlhf,agent,machine-learning}
\end{entry}

\begin{entry}{rlhf}{Reinforcement Learning from Human Feedback}[RLHF]
  \defn{christiano-2017}{Learning a reward function from human judgements
    comparing pairs of agent behaviours, and optimising a policy against that
    learned reward, where writing a reward by hand is impractical.}
  \defn{ouyang-2022}{A procedure for aligning language models with user
    intent. A pre-trained model is fine-tuned on demonstrations, a reward
    model is trained on human rankings of outputs, and the model is then
    optimised against the reward model with reinforcement learning.}
  \defn{ji-2023-alignment}{A family of \emph{learning from feedback}
    techniques used for forward alignment, in which human preference signals
    steer model behaviour toward intended values.}
  \related{alignment,fine-tuning,reinforcement-learning}
\end{entry}

\begin{entry}{self-supervised-learning}{Self-Supervised Learning}[SSL]
  \defn{liu-2021-ssl}{Learning in which supervisory signals are derived
    automatically from the data itself, for example by predicting masked or
    future parts of the input. It needs no human-annotated labels.}
  \defn{bommasani-2021}{The training method behind foundation models. A
    model is trained on unlabelled data with an objective derived from the
    data, such as predicting the next word.}
  \defn{iso-22989}{Machine learning that uses unlabelled data and
    automatically generated labels. It is sometimes treated as a subset of
    unsupervised learning.}
  \related{foundation-model,unsupervised-learning}
\end{entry}

\begin{entry}{supervised-learning}{Supervised Learning}
  \defn{iso-22989}{Machine learning that uses only labelled data during
    training.}
  \defn[ch.~2]{hastie-2009}{Learning to predict an outcome from inputs using a
    training set of observed input–outcome pairs, where the outcome acts as
    a \enquote{teacher} guiding learning.}
  \defn[§5.1.3]{goodfellow-2016}{Learning from a dataset in which each
    example is associated with a label or target, typically in order to
    estimate $p(y \mid x)$.}
  \related{unsupervised-learning,self-supervised-learning,datasets}
\end{entry}

\begin{entry}{transfer-learning}{Transfer Learning}
  \defn{pan-yang-2010}{Improving learning of a target task in a target
    domain by using knowledge from a different source domain or task,
    particularly when target data are scarce.}
  \defn[§15.2]{goodfellow-2016}{Using what has been learned in one setting,
    such as one distribution or task, to improve generalisation in another
    setting.}
  \defn{bommasani-2021}{The mechanism that makes foundation models possible:
    knowledge learned during pre-training on a surrogate task is transferred
    to downstream tasks through adaptation.}
  \related{fine-tuning,foundation-model}
\end{entry}

\begin{entry}{transformer}{Transformer}
  \defn*{vaswani-2017}{A new simple network architecture, the Transformer,
    based solely on attention mechanisms, dispensing with recurrence and
    convolutions entirely.}
  \defn{zhao-2023}{The standard backbone of large language models. It uses
    stacked multi-head self-attention and feed-forward layers, and its
    parallelism allows training on very large corpora.}
  \defn{bommasani-2021}{An architecture whose scalability and ability to
    handle multiple modalities made it the dominant basis of foundation
    models.}
  \related{large-language-model,neural-network,deep-learning}
\end{entry}

\begin{entry}{unsupervised-learning}{Unsupervised Learning}
  \defn{iso-22989}{Machine learning that uses only unlabelled data during
    training.}
  \defn[ch.~14]{hastie-2009}{Learning without a \enquote{teacher}: inferring
    properties of the probability density of the inputs directly, for
    example through clustering, density estimation, or dimensionality
    reduction.}
  \defn[§5.1.3]{goodfellow-2016}{Learning useful properties of the structure
    of a dataset, often the probability distribution that generated it,
    from examples without labels.}
  \related{supervised-learning,self-supervised-learning}
\end{entry}
